\item \textbf{Gradiente de índice de refracción en un tanque.}
Una solución de jarabe de maíz y agua se vierte en un recipiente de manera que su índice de refracción aumenta uniformemente desde $n_1 = 1.33$ en la superficie ($y = 20.0\text{ cm}$) hasta $n_2 = 1.90$ en el fondo ($y = 0$).
\begin{enumerate}
    \item Escriba la función del índice de refracción $n(y)$ a lo largo de la altura del tanque.
    \item Calcule el camino óptico total recorrido por un rayo vertical que cruza los $20.0\text{ cm}$ del tanque resolviendo:
    $$ \text{Camino Óptico} = \int_0^{20\text{ cm}} n(y) dy $$
    \item Si un haz de luz incide con un ángulo no nulo con respecto a la normal a la superficie de la mezcla, describa cualitativamente su trayectoria en el interior del tanque.
\end{enumerate}